% =============================================================
% Capitulo 9 - Desarrollo del software de interfaz (HMI)
% =============================================================
\chapter{Desarrollo del software de interfaz (HMI)}
\label{cap:software_hmi}

\section{Requisitos funcionales y de dise\~no}
\label{sec:requisitos_hmi}

El software HMI debe permitir la interacci\'on entre el usuario y el
microcontrolador de la placa de control, incluyendo la visualizaci\'on
del estado de sintonizaci\'on del sistema, la lectura de par\'ametros
del lazo de control y la modificaci\'on de funciones y \emph{setpoints}
de operaci\'on \cite{sat2026}. Su dise\~no toma como base los
lineamientos de la norma ISA-101 introducidos en la
Secci\'on~\ref{sec:isa101} \cite{informeAbril2026}.

\section{Arquitectura del software en Python}
\label{sec:arquitectura_python}

Se implement\'o un primer borrador del HMI en Python, utilizando la
biblioteca Tkinter para la interfaz gr\'afica y la biblioteca
\texttt{serial} para la comunicaci\'on por puerto USB con la placa de
control \cite{informeAbril2026}. El c\'odigo fuente del HMI se
documenta y versiona en el repositorio del proyecto \cite{kicad_hsrl}.

% PLACEHOLDER: describir la arquitectura definitiva del software (
% modulos, manejo de hilos/eventos para la comunicacion serie,
% estructura de pantallas) una vez avanzado su desarrollo.

\section{Comunicaci\'on con la placa de control}
\label{sec:comunicacion_placa}

Al momento de redacci\'on de este informe, se encuentra funcional
\'unicamente la comunicaci\'on en sentido placa-PC, permitiendo obtener
desde el HMI el estado de sintonizaci\'on del sistema y sus par\'ametros
\cite{informeAbril2026}.

PLACEHOLDER: completar esta secci\'on con la implementaci\'on de la
comunicaci\'on en sentido PC-placa (env\'io de comandos y
\emph{setpoints}) y con las capturas de pantalla del HMI en su versi\'on
final.
